\subsection{Integration on a fundamental domain}

Let X be a locally compact space, H a discrete group operating on the right continuously and properly in X. Let \(\pi\) be the canonical mapping of X onto X/H. For every \(x \in X\) , we denote by \(H_x\) the stabilizer of x in H; this is a finite subgroup of H (GT, III, §4, No.~2, Prop. 4); its order will be denoted \(n(x)\) . For every \(s \in H\) , \(H_{xs} = s^{-1}H_x s\) , therefore \(n(xs) = n(x)\) . There exists an open neighborhood U of x such that \(U \cap Us = \varnothing\) for \(s \notin H_x\) (loc. cit., No.~4, proof of Prop. 8); for \(y \in U\) , one has \(H_y \subset H_x\) ; thus the function n on X is upper semi-continuous. When X is countable at infinity, H is countable; for, let \((\mathrm{K}_1, \mathrm{K}_2, \ldots)\) be a covering of X by a sequence of compact subsets, and let \(x_{0} \in X\) ; the set of \(s \in H\) such that \(x_{0}s \in K_{i}\) is finite (loc. cit., No.~5, Th. 1), whence our assertion.

DEFINITION 2. --- Let \(F \subset X\) . One says that F is a fundamental domain (for H) if the restriction of \(\pi\) to F is a bijection of F onto X/H (in other words, F is a system of representatives for the equivalence relation defined by H).

Lemma 8. --- Let F be a fundamental domain. For every \(x \in \mathbf{X}\) ,

\[
\sum_ {s \in \mathrm{H}} \varphi_ {\mathrm{F} s} (x) = n (x).\tag{24}
\]

Since \(\varphi_{\mathrm{Fs}}(xt) = \varphi_{\mathrm{Fst}^{-1}}(x)\) for all \(s\) and \(t\) in \(\mathbf{H}\) , the two members of (24) remain invariant when \(x\) is replaced by \(xt\) . We can therefore suppose that \(x \in \mathbf{F}\) . We then have the equivalences

\[
\varphi_ {\mathrm{Fs}} (x) = 1 \Leftrightarrow x \in \mathrm{Fs} \Leftrightarrow x s ^ {- 1} \in \mathrm{F} \Leftrightarrow x s ^ {- 1} = x \Leftrightarrow s \in \mathrm{H} _ {x},
\]

whence (24).

PROPOSITION 15. --- Assume that X is countable at infinity. Let \(\mu\) be a measure \(\geqslant 0\) on X. Let F be a fundamental domain such that Fs is \(\mu\) -measurable for every \(s \in H\) . Let f be a \(\mu\) -integrable function on X, with values in a Banach space or in \(\overline{\mathbf{R}}\) . Then the family of the

\[
\int_ {\mathrm{Fs}} n (x) ^ {- 1} \mathbf {f} (x) d \mu (x) \quad (s \in \mathbf {H})
\]

is summable, and

\[
\int_ {\mathrm{X}} \mathbf {f} (x) d \mu (x) = \sum_ {s \in \mathrm{H}} \int_ {\mathrm{F} s} n (x) ^ {- 1} \mathbf {f} (x) d \mu (x).
\]

If \(\mathbf{A}\) is a finite subset of \(\mathbf{H}\) , then

\[
\left| \sum_ {s \in \mathrm{A}} n ^ {- 1} \mathbf {f} \varphi_ {\mathrm{Fs}} \right| \leqslant n ^ {- 1} | \mathbf {f} | \sum_ {s \in \mathrm{A}} \varphi_ {\mathrm{Fs}} \leqslant | \mathbf {f} |
\]

by Lemma 8. Lemma 8 also proves that \(\sum_{s\in\mathbf{A}}n^{-1}\mathbf{f}\varphi_{\mathrm{F}s}\) converges pointwise to \(\mathbf{f}\) with respect to the increasing directed set of finite subsets of \(\mathbf{H}\) . Prop. 15 then follows from Ch. IV, §4, No.~3, Th. 2.

THEOREM 4. --- Let X be a locally compact space countable at infinity, H a discrete group operating continuously and properly on the right in X, \(\pi\) the canonical mapping of X onto X/H, \(\mu\) a positive measure on X invariant under H, \(\beta\) the normalized Haar measure of H, and \(\lambda = \mu/\beta\) . Let F be a \(\mu\) -measurable fundamental domain.

\begin{enumerate}
\def\labelenumi{\alph{enumi})}
\tightlist
\item
  The pair \((\pi, n^{-1}\varphi_{\mathrm{F}})\) is \(\mu\) -adapted, and
\end{enumerate}

\[
\int_ {\mathrm{X}} n (x) ^ {- 1} \varphi_ {\mathrm{F}} (x) \varepsilon_ {\pi (x)} d \mu (x) = \lambda .
\]

\begin{enumerate}
\def\labelenumi{\alph{enumi})}
\setcounter{enumi}{1}
\item
  The mapping \(\pi\) is proper for \(n^{-1}\varphi_{\mathrm{F}}\cdot \mu\) , and \(\pi (n^{-1}\varphi_{\mathrm{F}}\cdot \mu) = \lambda\) .
\item
  Let k be a function on X/H. For k to be \(\lambda\) -measurable (resp. \(\lambda\) -integrable), it is necessary and sufficient that \(n^{-1}\varphi_{\mathrm{F}}(k\circ\pi)\) be \(\mu\) -measurable (resp. \(\mu\) -integrable); and, if k is \(\lambda\) -integrable then
\end{enumerate}

\[
\int_ {\mathrm{X/H}} k d \lambda = \int_ {\mathrm{F}} n ^ {- 1} (k \circ \pi) d \mu .
\]

We have \(\mu = \lambda^{\sharp}\) . Let \(f \in \mathcal{H}_{+}(\mathrm{X / H})\) . Then \(n^{-1}\varphi_{\mathrm{F}}(f \circ \pi)\) is \(\mu\) -measurable and \(\geqslant 0\) , and by Prop. 5 b) of No.~3 we have

\[
\int_ {\mathrm{X}} ^ {*} n (x) ^ {- 1} \varphi_ {\mathrm{F}} (x) f (\pi (x)) d \mu (x) = \int_ {\mathrm{X/H}} ^ {*} f (\dot {x}) d \lambda (\dot {x}) \int_ {\mathrm{H}} ^ {*} n (x \xi) ^ {- 1} \varphi_ {\mathrm{F}} (x \xi) d \beta (\xi)
\]

and \(\int_{\mathrm{H}}^{*}n(x\xi)^{-1}\varphi_{\mathrm{F}}(x\xi)d\beta (\xi) = n(x)^{-1}\sum_{\xi \in \mathrm{H}}\varphi_{\mathrm{F}}(x\xi) = 1\) by Lemma 8. Thus \(n^{-1}\varphi_{\mathrm{F}}\cdot (f\circ \pi)\) is \(\mu\) -integrable and

\[
\int_ {\mathrm{X}} n (x) ^ {- 1} \varphi_ {\mathrm{F}} (x) f (\pi (x)) d \mu (x) = \int_ {\mathrm{X/H}} f (\dot {x}) d \lambda (\dot {x}).
\]

This proves a). The assertion b) is proved similarly. The assertion c) may be deduced from a) and from Ch. V, §4, Prop. 3 and Th. 2.

COROLLARY. --- We maintain the hypotheses and notations of Th. 4. Let F' be a second μ-measurable fundamental domain. Let u be a function on X, with values in a Banach space or in \(\overline{R}\) , invariant under H. Suppose that u is μ-integrable on F. Then u is μ-integrable on F' and

\[
\int_ {\mathrm{F}} \mathbf {u} (x) d \mu (x) = \int_ {\mathrm{F} ^ {\prime}} \mathbf {u} (x) d \mu (x).
\]

Since \(\mathbf{u}\) and \(n\) are invariant under \(\mathbf{H}\) , there exists a function \(\mathbf{v}\) on \(\mathrm{X / H}\) such that \(\mathbf{v} \circ \pi\) coincides with \(n\mathbf{u}\) on \(\mathbf{F}\) and on \(\mathbf{F}'\) . Then \(n^{-1}\varphi_{\mathrm{F}}(\mathbf{v}\circ \pi) = \varphi_{\mathrm{F}}\mathbf{u}, n^{-1}\varphi_{\mathrm{F}^{\prime}}(\mathbf{v}\circ \pi) = \varphi_{\mathrm{F}^{\prime}}\mathbf{u}.\) By hypothesis, \(n^{-1}\varphi_{\mathrm{F}}(\mathbf{v}\circ \pi)\) is \(\mu\) -integrable. By Th. 4, \(\mathbf{v}\) is \(\lambda\) -integrable, \(\varphi_{\mathrm{F}^{\prime}}\mu\) is \(\mu\) -integrable, and

\[
\int_ {\mathrm{F}} \mathbf {u} d \mu = \int_ {\mathrm{X/H}} \mathbf {v} d \lambda = \int_ {\mathrm{F} ^ {\prime}} \mathbf {u} d \mu .
\]

For the existence of \(\mu\) -measurable fundamental domains, see Exer. 12.

